\documentclass[11pt,a4paper]{article}
\usepackage[utf8]{inputenc}
\usepackage[ngerman]{babel}
\usepackage{amsmath}
\usepackage{geometry}
\usepackage{parskip}
\geometry{a4paper, left=2.5cm, right=2.5cm, top=2.5cm, bottom=2.5cm}

\title{\textbf{Prüfungsvorbereitung Energy Systems} \\ \Large Die 60 wichtigsten Verständnisfragen}
\author{}
\date{}

\begin{document}

\maketitle

\section*{Global Energy Resources}

\textbf{1. The energy flow from mining to consumption is categorized in five stages. List these stages.} [cite: 2]\\
\textit{Antwort:} Primary Energy, Transformation Process, Secondary Energy, Final Energy Consumption, Useful Energy. [cite: 3]

\vspace{0.5cm}
\textbf{2. Define the meaning of these two terms Reserve and Resource.} [cite: 7]\\
\textit{Antwort:} Reserve: Mit heutiger Technologie wirtschaftlich förderbar. Ressource: Potenziell vorhanden, aber derzeit nicht wirtschaftlich abbaubar. [cite: 8]

\vspace{0.5cm}
\textbf{3. Name three non-renewable energy sources.} [cite: 12]\\
\textit{Antwort:} Fossile Brennstoffe (Kohle, Gas, Öl), Uran. [cite: 12]

\section*{Energy Transformation}

\textbf{4. Name and describe the levels of a traditional hierarchical power system structure.} [cite: 17]\\
\textit{Antwort:} Produktion (HV) $\rightarrow$ Übertragung (HV) $\rightarrow$ Verteilung (MV/LV) an den Endverbraucher. [cite: 18]

\vspace{0.5cm}
\textbf{5. Why does the hierarchical power system has several voltage levels?} [cite: 19]\\
\textit{Antwort:} Um große Energiemengen mit minimalen Leistungsverlusten über große Distanzen zu transportieren. [cite: 20]

\section*{Electrical Energy Production}

\textbf{6. Describe the terms base load, mid-merit load and peak load.} [cite: 26]\\
\textit{Antwort:} Grundlast (dauerhafter Bedarf), Mittellast (variabler Bedarf), Spitzenlast (kurzzeitige Leistungsspitzen). [cite: 27]

\vspace{0.5cm}
\textbf{7. What are typical base load power plants and what is their common characteristic?} [cite: 28]\\
\textit{Antwort:} Laufwasserkraftwerke, Nuklear- und Braunkohlekraftwerke. Charakteristik: Laufen kontinuierlich, schwer regelbar. [cite: 29]

\vspace{0.5cm}
\textbf{8. What are typical peak power plants and what is their common characteristic?} [cite: 32]\\
\textit{Antwort:} Pumpspeicherkraftwerke und Gasturbinen. Charakteristik: Sehr schnelle Startzeiten. [cite: 32]

\vspace{0.5cm}
\textbf{9. Describe the term spinning reserve and what is it used for?} [cite: 33]\\
\textit{Antwort:} Sofort verfügbare Erzeugungskapazität durch Generatoren, die synchron am Netz sind. Dient dem Ausgleich plötzlicher Lastwechsel oder Ausfälle. [cite: 34]

\vspace{0.5cm}
\textbf{10. What are the biggest differences between regular economic market and the electricity market?} [cite: 43]\\
\textit{Antwort:} Elektrische Energie ist in großen Mengen kaum speicherbar, Erzeugung Verbrauch müssen zeitgleich sein und der Stromfluss folgt physikalischen Gesetzen. [cite: 44]

\vspace{0.5cm}
\textbf{11. Why is the tradable product Electrical Energy time based?} [cite: 47]\\
\textit{Antwort:} Wegen der fehlenden großflächigen Speicherbarkeit muss Erzeugung und Verbrauch zeitgleich erfolgen. [cite: 48]

\section*{Power Generation Costs \& Planning}

\textbf{12. What are CAPEX and what are they used for?} [cite: 57]\\
\textit{Antwort:} Capital Expenditures (Investitionskosten). Decken den Bau und die Errichtung ab. [cite: 58]

\vspace{0.5cm}
\textbf{13. What are OPEX and what are they used for?} [cite: 61]\\
\textit{Antwort:} Operational Expenditures (Betriebskosten). Fallen für den laufenden Betrieb (z.B. Brennstoffe, Wartung) an. [cite: 62]

\vspace{0.5cm}
\textbf{14. How is the power generation costs defined?} [cite: 67]\\
\textit{Antwort:} Jährliche Gesamtkosten geteilt durch die jährlich produzierte elektrische Energie. [cite: 68]

\vspace{0.5cm}
\textbf{15. What is the accumulated daily load curve used for?} [cite: 70]\\
\textit{Antwort:} Um die Lastsituation zu definieren und die richtigen Kraftwerkstechnologien sowie deren minimale Betriebsstunden zu planen. [cite: 71]

\section*{Energy Trading}

\textbf{16. What markets are available for electrical power?} [cite: 78]\\
\textit{Antwort:} Spotmärkte (Day-Ahead, Intraday) und Terminmärkte (Futures). [cite: 79]

\vspace{0.5cm}
\textbf{17. What are the two basic ways business transactions are happening?} [cite: 86]\\
\textit{Antwort:} Kontinuierlicher Handel (Continuous trading) und Auktionen. [cite: 87]

\vspace{0.5cm}
\textbf{18. Which type of transaction is using an open ordering book and how does it works?} [cite: 88]\\
\textit{Antwort:} Kontinuierlicher Handel. Kauf- und Verkaufsaufträge werden gesammelt und fortlaufend abgeglichen. [cite: 89]

\vspace{0.5cm}
\textbf{19. Which type of transaction is using a closed ordering book and how does it works?} [cite: 92]\\
\textit{Antwort:} Auktionen. Aufträge werden bis zu einer Deadline gesammelt, danach wird ein einheitlicher Preis gebildet. [cite: 93]

\vspace{0.5cm}
\textbf{20. What is a Market Clearing Price?} [cite: 96]\\
\textit{Antwort:} Der Schnittpunkt der aggregierten Angebots- und Nachfragekurven; ein für alle Teilnehmer gültiger Preis. [cite: 97]

\vspace{0.5cm}
\textbf{21. What does a physical fulfillment of a business transaction means and for what electricity products is it usually used?} [cite: 100]\\
\textit{Antwort:} Tatsächliche Stromlieferung gegen Geld. Genutzt bei Spotmärkten (Day-Ahead/Intraday) und physischen Optionen. [cite: 101]

\vspace{0.5cm}
\textbf{22. What does the abbreviation OTC trading stays for and what does it means?} [cite: 108]\\
\textit{Antwort:} Over-the-counter. Direkter bilateraler Handel außerhalb der Strombörse. [cite: 109]

\vspace{0.5cm}
\textbf{23. How does the Day Ahead Trading system works?} [cite: 119]\\
\textit{Antwort:} Aufträge werden bis 12:00 Uhr des Vortages gesammelt. Per Auktion wird der Market Clearing Price für jede Stunde des Folgetages ermittelt. [cite: 120]

\vspace{0.5cm}
\textbf{24. What is the main decision factor when a utility has to decide to produce electrical energy by them self or buying it from the market...} [cite: 121]\\
\textit{Antwort:} Der Market Clearing Price im Vergleich zu den eigenen variablen Produktionskosten (Grenzkosten). [cite: 122]

\vspace{0.5cm}
\textbf{25. What does the Merit order curve shows?} [cite: 123]\\
\textit{Antwort:} Einsatzreihenfolge der Kraftwerke, sortiert nach ihren Grenzkosten (von günstig zu teuer). [cite: 124]

\vspace{0.5cm}
\textbf{26. How does the Intraday Trading system works?} [cite: 130]\\
\textit{Antwort:} Fortlaufender Handel (offenes Orderbuch) bis kurz vor Lieferbeginn, um kurzfristige Abweichungen auszugleichen. [cite: 131]

\vspace{0.5cm}
\textbf{27. What is the Futures Trading system and what is it used for?} [cite: 132]\\
\textit{Antwort:} Handel von Strommengen zu einem Fixpreis weit in der Zukunft, zur langfristigen Preisabsicherung (Hedging). [cite: 133]

\vspace{0.5cm}
\textbf{28. What is a balance group used for and who can participate in such a group?} [cite: 136]\\
\textit{Antwort:} Bündelt Netznutzer zum Bilanzabgleich zwischen Ein- und Ausspeisung. Jeder Marktteilnehmer muss Mitglied einer Bilanzgruppe sein. [cite: 137]

\section*{Thermal Plants}

\textbf{29. What is a cycle process (dt. Kreiprozess)?} [cite: 143]\\
\textit{Antwort:} Ein thermodynamischer Vorgang, bei dem das Arbeitsmedium nach mehreren Zustandsänderungen wieder in den Ausgangszustand zurückkehrt. [cite: 144]

\vspace{0.5cm}
\textbf{30. Describe the four stages of the Carnot cycle.} [cite: 147]\\
\textit{Antwort:} 1-2 Isotherme Expansion (Wärmezufuhr), 2-3 Isentrope Expansion, 3-4 Isotherme Kompression (Wärmeabfuhr), 4-1 Isentrope Kompression. [cite: 148]

\vspace{0.5cm}
\textbf{31. Describe the four stages of the Rankine cycle.} [cite: 155]\\
\textit{Antwort:} 1-2 Isentrope Kompression (Pumpe), 2-3 Isobare Wärmezufuhr/Verdampfung, 3-4 Isentrope Entspannung (Turbine), 4-1 Isobare Kondensation. [cite: 156]

\vspace{0.5cm}
\textbf{32. What is the difference between an ideal and an real Rankine process.} [cite: 159]\\
\textit{Antwort:} Der reale Prozess beinhaltet Verluste: Druckverluste in Rohren und eine nicht-isentrope (verlustbehaftete) Expansion in der Turbine. [cite: 160]

\vspace{0.5cm}
\textbf{33. Describe the four stages of the Brayton cycle.} [cite: 172]\\
\textit{Antwort:} 1-2 Isentrope Kompression, 2-3 Isobare Wärmezufuhr, 3-4 Isentrope Entspannung, 4-1 Isobare Wärmeabfuhr. [cite: 173]

\vspace{0.5cm}
\textbf{34. Name the components of a coal fired power plant.} [cite: 176]\\
\textit{Antwort:} Kohlemühle, Kessel, Dampfturbine, Kondensator, Speisewasserpumpe, Elektrofilter, Rauchgasentschwefelungsanlage (REA), DeNOx-Anlage. [cite: 177]

\vspace{0.5cm}
\textbf{35. What are the two basic different designs of a steam generator?} [cite: 194]\\
\textit{Antwort:} Naturumlaufkessel und Zwangsumlaufkessel (bzw. Zwangdurchlaufkessel/Benson). [cite: 195]

\vspace{0.5cm}
\textbf{36. What are condenser used for in a thermal power plant?} [cite: 198]\\
\textit{Antwort:} Um den aus der Turbine austretenden Dampf wieder zu flüssigem Wasser zu kondensieren und so den Kreislauf zu schließen. [cite: 199]

\vspace{0.5cm}
\textbf{37. Describe how a wet scrubber is used for flue gas desulphurization.} [cite: 204]\\
\textit{Antwort:} Rauchgas wird mit einer Kalkstein-Wasser-Suspension besprüht. Das $SO_{2}$ reagiert mit dem Kalk und fällt als Gips aus. [cite: 205]

\vspace{0.5cm}
\textbf{38. Describe the Reaction Principle of a steam turbine.} [cite: 212]\\
\textit{Antwort:} Der Dampfdruck fällt sowohl im Leitrad als auch im Laufrad ab. Die Entspannung im Laufrad erzeugt eine Reaktionskraft, die den Rotor antreibt. [cite: 213, 214]

\vspace{0.5cm}
\textbf{39. Describe the Impulse Principle of a steam turbine.} [cite: 215]\\
\textit{Antwort:} Der gesamte Druckabbau erfolgt im Leitrad (hohe Geschwindigkeit). Im Laufrad wird nur die kinetische Energie (Impuls) ohne weiteren Druckabfall übertragen. [cite: 216]

\vspace{0.5cm}
\textbf{40. What is the working principle of combined-cycle gas-turbine (CCGT) plants?} [cite: 217]\\
\textit{Antwort:} Gas- und Dampfturbinenkraftwerk (GuD): Die heißen Abgase der Gasturbine erzeugen in einem Abhitzekessel Dampf, der eine Dampfturbine antreibt. [cite: 218]

\vspace{0.5cm}
\textbf{41. Describe boiling water reactor (BWR) and Pressurized Water Reactor (PWR).} [cite: 227]\\
\textit{Antwort:} BWR: Wasser siedet direkt im Reaktorkern; Dampf treibt die Turbine. [cite: 228] PWR: Wasser im Primärkreis steht unter hohem Druck und siedet nicht; ein Wärmetauscher erzeugt Dampf im Sekundärkreis. [cite: 229]

\vspace{0.5cm}
\textbf{42. Describe the control strategy Boiler following.} [cite: 232]\\
\textit{Antwort:} Die Turbine regelt die elektrische Leistung. Fällt der Dampfdruck dadurch ab, regelt der Kessel die Feuerung nach, um den Druck wiederherzustellen. [cite: 233]

\vspace{0.5cm}
\textbf{43. Describe the control strategy Turbine following.} [cite: 234]\\
\textit{Antwort:} Der Kessel wird auf eine bestimmte Leistung geregelt. Die Turbinenventile reagieren auf den entstehenden Dampfdruck und öffnen/schließen entsprechend. [cite: 235]

\section*{Renewable Plants}

\textbf{44. What is the commonly used turbine for river power plants, and why?} [cite: 245]\\
\textit{Antwort:} Kaplan-Turbine, da sie bei geringen Fallhöhen und stark schwankenden Wassermengen den höchsten Wirkungsgrad erzielt. [cite: 246]

\vspace{0.5cm}
\textbf{45. What is a commonly used turbine used for storage power plants, and why?} [cite: 249]\\
\textit{Antwort:} Francis-Turbinen (mittlere bis hohe Fallhöhen) und Pelton-Turbinen (sehr hohe Fallhöhen), da sie für diese Drücke optimal konstruiert sind. [cite: 250]

\vspace{0.5cm}
\textbf{46. Describe a typical cyclic operation (dt. Wälzbetrieb) of an pump storage power plant.} [cite: 255]\\
\textit{Antwort:} Tagsüber (Spitzenlast/hohe Preise) Wasser abarbeiten und Strom erzeugen; nachts (Schwachlast/niedrige Preise) Wasser wieder ins Oberbecken pumpen. [cite: 256]

\vspace{0.5cm}
\textbf{47. Describe what the operation mode hydraulic short circuit means and what it is used for.} [cite: 263]\\
\textit{Antwort:} Pumpe und Turbine laufen gleichzeitig. Da große Pumpen oft nicht regelbar sind, steuert die Turbine den Bypass, um die Netto-Aufnahmeleistung stufenlos und hochdynamisch fürs Stromnetz (Frequenzregelung) zu justieren. [cite: 264]

\vspace{0.5cm}
\textbf{48. Describe the phenomenon cavitation erosion.} [cite: 273]\\
\textit{Antwort:} Bei lokalem Unterdruck verdampft Wasser. Steigt der Druck wieder, implodieren diese Dampfblasen schlagartig und reißen Material aus den Turbinenschaufeln. [cite: 274]

\vspace{0.5cm}
\textbf{49. Describe how a Francis turbine works.} [cite: 278]\\
\textit{Antwort:} Reaktionsturbine: Wasser strömt radial in das spiralförmige Gehäuse ein, durchströmt das Laufrad und tritt axial nach unten aus. Druck und Geschwindigkeit erzeugen Drehmoment. [cite: 279, 280]

\vspace{0.5cm}
\textbf{50. Describe how a Pelton turbine works.} [cite: 281]\\
\textit{Antwort:} Aktions-/Gleichdruckturbine: Hochgeschwindigkeits-Wasserstrahlen treffen auf becherförmige Schaufeln des Laufrades und übertragen ihre kinetische Energie durch Impulsumkehr. [cite: 282]

\vspace{0.5cm}
\textbf{51. Describe how a Kaplan turbine works.} [cite: 283]\\
\textit{Antwort:} Wie ein Schiffspropeller. Sowohl Leitapparat als auch Laufradschaufeln sind verstellbar, um sich optimal an wechselnde Wasserführungen anzupassen. [cite: 284]

\vspace{0.5cm}
\textbf{52. Describe how photovoltaic conversion works.} [cite: 285]\\
\textit{Antwort:} Sonnenlicht (Photonen) regt Elektronen in einem dotierten Halbleiter (Silizium) an. Das p-n-Feld trennt die Ladungen, wodurch eine elektrische Gleichspannung entsteht. [cite: 286]

\vspace{0.5cm}
\textbf{53. What makes the blades of a wind turbine move? Describe the physical principle.} [cite: 289]\\
\textit{Antwort:} Der aerodynamische Auftrieb: Der Wind umströmt das asymmetrische Profil, erzeugt Unterdruck auf der Oberseite und zieht das Blatt nach vorn. [cite: 290]

\vspace{0.5cm}
\textbf{54. Describe the working principle of an doubly fed induction generator (dt. DoppelGespeisten Asynchrongenerator).} [cite: 293]\\
\textit{Antwort:} Der Stator ist direkt mit dem Netz verbunden. Der Rotor wird über einen Frequenzumrichter gespeist, wodurch die Anlage mit variabler Drehzahl betrieben werden kann, bei konstanter Netzfrequenz. [cite: 294]

\vspace{0.5cm}
\textbf{55. Name the two ways biomass can be used to produce electricity.} [cite: 303]\\
\textit{Antwort:} 1. Direkte Verbrennung (z.B. Holz) in einem Kessel zur Dampferzeugung für eine Dampfturbine. [cite: 304] 2. Vergärung zu Biogas, das in einem Verbrennungsmotor/Gasturbine verstromt wird. [cite: 305]

\section*{Power System Control}

\textbf{56. What are the tasks of a control area manager (dt. Regelzonenführer)?} [cite: 307]\\
\textit{Antwort:} Sicherstellung der Systembilanz (Erzeugung gleich Verbrauch), Frequenzhaltung, Beschaffung von Regelleistung und Netzengpassmanagement. [cite: 308]

\vspace{0.5cm}
\textbf{57. Describe the working principle of an synchronous generator?} [cite: 309]\\
\textit{Antwort:} Ein rotierendes, gleichstromerregtes Magnetfeld im Rotor induziert in den feststehenden Wicklungen des Stators eine Wechselspannung (Netzfrequenz ist starr an Drehzahl gekoppelt). [cite: 310]

\vspace{0.5cm}
\textbf{58. What are the three-stage procedure consists of, which is used to keep the electrical power balanced within the Continental European Grid?} [cite: 315]\\
\textit{Antwort:} Primärregelung, Sekundärregelung, Tertiärregelung (Minutenreserve). [cite: 316]

\vspace{0.5cm}
\textbf{59. What is the time period for a single incident in the primary control stage?} [cite: 317]\\
\textit{Antwort:} Reaktion erfolgt innerhalb von Sekunden und muss bis zu 30 Sekunden vollständig erbracht sein. [cite: 318]

\vspace{0.5cm}
\textbf{60. What is the time period for a single incident in the secondary control stage?} [cite: 321]\\
\textit{Antwort:} Aktivierung ab ca. 30 Sekunden, muss innerhalb von 5 bis 15 Minuten vollständig erbracht sein. [cite: 322]

\end{document}